\section{Normalized congruence triviality and the level-three gate}
\label{sec:normalized-congruence}

The level-two identity has a consequence at every congruence level.
Use the signed representation
$\widetilde q:\widehat B_3\to\operatorname{SL}_2(\widehat{\mathbb Z})$
with $u\mapsto U=\left(\begin{smallmatrix}1&1\\0&1\end{smallmatrix}\right)$
and $v\mapsto V=\left(\begin{smallmatrix}1&0\\-1&1\end{smallmatrix}\right)$.

\begin{proposition}\label{prop:normalized-congruence-triviality}
Let $g=(1,f)\in\widehat{\mathrm{GT}}$ satisfy $\rho_2(g)=0$.
Then
\[
 f(U^2,V^2)=I
       \qquad\text{in }\operatorname{SL}_2(\widehat{\mathbb Z}).
\]
Consequently $\widetilde q\Phi_g=\widetilde q$ on $\widehat B_3$.
\end{proposition}
\begin{proof}
Normalized relation (IV), together with interchange of the two braid
generators and the two-cycle relation, gives
\[
 f(u^4,v)=f(u^2,v^2)
       \qquad\text{in }\widehat B_3.
\]
Fix a prime $\ell$ and put $Q=f(U^2,V^2)$ in
$\operatorname{SL}_2(\mathbb Z_\ell)$.
Conjugation by $D=\operatorname{diag}(2,1)$ in
$\operatorname{GL}_2(\mathbb Q_\ell)$ sends $U^2$ to $U^4$ and
$V^2$ to $V$. Therefore the displayed signed identity gives
$DQD^{-1}=Q$, with no scalar ambiguity. Since the two diagonal
entries of $D$ are distinct, $Q=\operatorname{diag}(a,a^{-1})$
for a unit $a\in\mathbb Z_\ell^\times$.

The images of the braid generators under
$\widetilde q_\ell\Phi_g$ are $U$ and
\[
 M=Q^{-1}VQ=\begin{pmatrix}1&0\\-a^2&1\end{pmatrix}.
\]
Their braid relation $UMU=MUM$, in its lower-left entry, is
$\gamma(1+\gamma)=0$ with $\gamma=-a^2$.
The coefficient ring is a domain and $a$ is a unit, so $a^2=1$.
Thus $Q=\pm I$ and $M=V$.

It remains to determine the sign using the three-cycle relation.
Put $A=U^2$, $B=V^2$, $Z=(AB)^{-1}$, and $R=UV$.
Direct multiplication gives
\[
 RAR^{-1}=B,\qquad RBR^{-1}=-Z,\qquad RZR^{-1}=-A.
\]
A commutator word is unchanged when one of its arguments is multiplied
by a central scalar. Since $Q$ is central, conjugating its defining
expression by $R$ and $R^2$ consequently gives
\[
 f(A,B)=f(B,Z)=f(Z,A)=Q.
\]
The multiplier-one three-cycle relation, with $ABZ=1$, now implies
$Q^3=I$. Together with $Q^2=I$, this forces $Q=I$.
This holds for every $\ell$, proving the signed adelic assertion.
Both braid generators now have unchanged images, which proves the
intertwining identity.
\end{proof}

In the following statements the ambient modular group is the
centerless group $B=\operatorname{PSL}_2(\mathbb Z)$ and its full
profinite completion, rather than the central braid extension. Put
\[
 \mathcal C=\ker\bigl(\widehat B\longrightarrow
      \operatorname{SL}_2(\widehat{\mathbb Z})/\{\pm I\}\bigr),
\]
where the scalar sign is the common sign in all components.

\begin{lemma}\label{lem:parabolic-centralizer}
The centralizer of $v$ in $\widehat B$ is
$\overline{\langle v\rangle}\simeq\widehat{\mathbb Z}$, and
it has trivial intersection with $\mathcal C$.
\end{lemma}
\begin{proof}
The free pure subgroup is
$\widehat F_2=\overline{\langle x\rangle}\amalg
\overline{\langle y\rangle}$, where $x=u^2,y=v^2$.
The centralizer theorem for a nonidentity element of a factor in a
free profinite product gives
$C_{\widehat F_2}(y)=\overline{\langle y\rangle}$
\cite[Theorem B, p.~155; Theorem B$'$, \S4, p.~159]{HerfortRibes1985}.
This theorem has no torsion assumption on the element of the factor.
If $c\in C_{\widehat B}(v)$, its image in $S_3$ centralizes the
transposition represented by $v$. Thus $cv^{-e}$ is pure for
some $e\in\{0,1\}$ and centralizes $y$. It lies in
$\overline{\langle y\rangle}$, proving the centralizer assertion.
Finally $v^t$ has congruence image represented by
$\left(\begin{smallmatrix}1&0\\-t&1\end{smallmatrix}\right)$.
It is a scalar sign only if $t=0$. This proves injectivity of the
procyclic subgroup in the congruence quotient and the last assertion.
\end{proof}

\begin{corollary}\label{cor:level-three-exact-gate}
For a normalized $g=(1,f)\in\widehat{\mathrm{GT}}$, put
$q=f(u^2,v^2)\in\widehat B$. The element $q$ lies in
$\mathcal C$, hence in $\overline{\Gamma_0(3)}$.
The automorphism $\Phi_g$ preserves that subgroup. For
$d_3:\overline{\Gamma_0(3)}\to\widehat B$ induced by
conjugation by $\operatorname{diag}(3,1)$, set
\[
 h_3=d_3(q)q^{-1}.
\]
Then $h_3\in\mathcal C$, and strict compatibility with both
level-three degeneracy maps is equivalent to the exact identity
\[
 d_3(q)=q.
\]
\end{corollary}
\begin{proof}
The preceding proposition shows that $q$ and the action of $g$
are trivial at every congruence level. Therefore $q$ belongs to
$\mathcal C$ and $\Phi_g$ preserves the level-three subgroup.
On its congruence image, $d_3$ is conjugation by the same diagonal
matrix: the equality holds on the discrete subgroup and extends
by continuity. Consequently $d_3(q)$ also has trivial congruence
image, and $h_3\in\mathcal C$.

The subgroup $\Gamma_0(3)$ is generated by $u,v^3$, and
$d_3(u)=u^3$, $d_3(v^3)=v$. The two proposed compositions agree
on $u$. On $v^3$ they agree exactly when
$h_3\in C_{\widehat B}(v)$. The preceding lemma and
$h_3\in\mathcal C$ make this equivalent to $h_3=1$.
\end{proof}

The corollary identifies the remaining full profinite identity. It does
not prove it: every congruence representation already kills $h_3$, so
congruence matrices alone cannot decide whether this element is trivial
in $\widehat B$.
